\begin{figure}[htbp]
	\centering
	\resizebox{\textwidth}{!}{%
		\begin{tikzpicture}[
			trave/.style={line width=2.2pt, line cap=round,
				line join=round, black},
			terra/.style={line width=0.7pt, black},
			nodoV/.style={circle, fill=white, draw=black,
				inner sep=0pt, minimum size=5pt, line width=0.8pt},
			carico/.style={->, >=stealth, line width=1.0pt, black},
			reaz/.style={->, >=stealth, line width=1.0pt, gray!60},
			forza/.style={->, >=stealth, line width=1.0pt, black},
			blob/.style={draw=gray!45, fill=gray!10, line width=0.7pt,
				rounded corners=14pt},
			]
			
			\def\L{2.8}
			\def\sc{0.72}    % lunghezza frecce forze concentrate
			\def\hLoad{0.65} % altezza blocco frecce carichi distribuiti
			\def\hTri{0.82}  % lunghezza max frecce carico triangolare
			
			%% ════════════════════════════════════════════════
			%% SEMIARCO SINISTRO  A–C–D
			%% ════════════════════════════════════════════════
			\begin{scope}[xshift=0cm, yshift=0cm]
				\coordinate (A)    at (0,   0   );
				\coordinate (C)    at (0,   \L  );
				\coordinate (Dsin) at (\L,  \L  );
				
				%% Blob
				\draw[blob] (-1.,-0.6) rectangle (\L+0.8, \L+1.10);
				
				%% Travi
				\draw[trave] (A) -- (C) -- (Dsin);
				
				%% ── Cerniera esterna in A ────────────────────
				\draw[terra] (A)--++(-0.22,-0.38)--++(0.44,0)--cycle;
				\draw[terra] (-0.30,-0.38)--(0.30,-0.38);
				\foreach \x in {-0.24,-0.12,0,0.12,0.24}
				\draw[very thin,black](\x,-0.38)--++(-0.09,-0.11);
				\node[nodoV] at (A) {};
				
				%% ── Cerniera interna in D ────────────────────
				\node[nodoV] at (Dsin) {};
				
				%% ── Reazioni in A ────────────────────────────
				%% H_A = (2/3)pℓ  verso destra (→)
				\draw[reaz] (A) -- ++(\sc, 0)
				node[above, font=\small, gray!70]
				{$\tfrac{2}{3}p\ell$};
				%% V_A = (4/3)pℓ  verso l'alto (↑)
				\draw[reaz] (A) -- ++(0, \sc)
				node[right, font=\small, gray!70]
				{$\tfrac{4}{3}p\ell$};
				
				%% ── Carico triangolare su AC (→) ─────────────
				%% Profilo triangolare
				\draw[black, line width=0.6pt]
				(A) -- ($(C)+(-\hTri, 0)$);
				%% Frecce (t=0 escluso: freccia di lunghezza nulla)
				\foreach \t in {0.15,0.30,0.45,0.60,0.75,0.90,1.00}{
					\pgfmathsetmacro{\lungh}{\t*\hTri}
					\draw[carico]
					($(A)!\t!(C)+(-\lungh,0)$)
					-- ($(A)!\t!(C)$);
				}
				\node[left, font=\small]
				at ($(C)+(-\hTri+0.3,0.2)$) {$p$};
				
				%% ── Carico uniforme su CD (↓) ────────────────
				\draw[black, line width=0.6pt]
				($(C)+(0,\hLoad)$) -- ($(Dsin)+(0,\hLoad)$);
				\draw[black, line width=0.4pt]
				($(C)+(0,0.07)$)   -- ($(Dsin)+(0,0.07)$);
				\foreach \t in {0.00,0.17,0.33,0.50,0.67,0.83,1.00}{
					\draw[carico, line width=0.8pt]
					($(C)!\t!(Dsin)+(0,\hLoad)$)
					-- ($(C)!\t!(Dsin)+(0,0.07)$);
				}
				\node[above, font=\small]
				at ($(C)!0.5!(Dsin)+(0,\hLoad)$) {$p$};
				
				%% ── Forze interne in D su semiarco sinistro ──
				%% H_D = (7/6)pℓ  verso sinistra (←)
				\draw[forza] ($(Dsin)+(\sc,0)$) -- (Dsin);
				\node[below, font=\small] at ($(Dsin)+(0.5*\sc,0)$) {$\tfrac{7}{6}p\ell$};
				
				%%				\draw[forza] (Dsin) -- ++(-\sc, 0)
				%%				node[below, font=\small]
				%%				{$\tfrac{7}{6}p\ell$};
				%% V_D = (1/3)pℓ  verso il basso (↓)
				\draw[forza] (Dsin) -- ++(0, -\sc)
				node[left, font=\small]
				{$\tfrac{1}{3}p\ell$};
				
				%% ── Etichette nodi ───────────────────────────
				\node[left,  font=\small] at (A)    {$A$};
				\node[above left,  font=\small]
				at ($(C)+(-0.05,0.08)$) {$C$};
				\node[above right, font=\small] at (Dsin) {$D$};
				
			\end{scope}
			
			
			%% ════════════════════════════════════════════════
			%% SEMIARCO DESTRO  D–E–B
			%% ════════════════════════════════════════════════
			\pgfmathsetmacro{\XD}{\L + 1.90}
			\begin{scope}[xshift=\XD cm, yshift=0cm]
				\coordinate (Ddes) at (0,   \L      );
				\coordinate (E)    at (\L,  \L      );
				\coordinate (B)    at (\L,  0       );
				\coordinate (F)    at (\L,  0.5*\L  );
				
				%% Blob
				\draw[blob] (-0.8,-0.6) rectangle (\L+1.0, \L+1.10);
				
				
				%% Travi
				\draw[trave] (Ddes) -- (E) -- (B);
				
				%% ── Cerniera esterna in B ────────────────────
				\draw[terra] (B)--++(-0.22,-0.38)--++(0.44,0)--cycle;
				\draw[terra] ($(B)+(-0.30,-0.38)$)
				-- ($(B)+(0.30,-0.38)$);
				\foreach \x in {-0.24,-0.12,0,0.12,0.24}
				\draw[very thin,black]
				($(B)+(\x,-0.38)$)--++(-0.09,-0.11);
				\node[nodoV] at (B) {};
				
				%% ── Cerniera interna in D ────────────────────
				\node[nodoV] at (Ddes) {};
				
				%% ── Punto F (trattino) ───────────────────────
				\draw[gray!60, line width=0.8pt]
				($(F)+(-0.14,0)$) -- ($(F)+(0.14,0)$);
				
				%% ── Forze interne in D su semiarco destro ────
				%% H_D = (7/6)pℓ  verso destra (→)  [azione e reazione]
				\draw[forza] ($(Ddes)+(-\sc,0)$) -- (Ddes);
				\node[below, font=\small] at ($(Ddes)+(-0.5*\sc,0)$) {$\tfrac{7}{6}p\ell$};
				%% V_D = (1/3)pℓ  verso l'alto (↑)
				\draw[forza] ($(Ddes)+(0,-\sc)$) -- (Ddes);
				\node[right, font=\small] at ($(Ddes)+(0,-0.5*\sc)$) {$\tfrac{1}{3}p\ell$};
				%% ── Carico uniforme su DE (↓) ────────────────
				\draw[black, line width=0.6pt]
				($(Ddes)+(0,\hLoad)$) -- ($(E)+(0,\hLoad)$);
				\draw[black, line width=0.4pt]
				($(Ddes)+(0,0.07)$)   -- ($(E)+(0,0.07)$);
				\foreach \t in {0.00,0.17,0.33,0.50,0.67,0.83,1.00}{
					\draw[carico, line width=0.8pt]
					($(Ddes)!\t!(E)+(0,\hLoad)$)
					-- ($(Ddes)!\t!(E)+(0,0.07)$);
				}
				\node[above, font=\small]
				at ($(Ddes)!0.5!(E)+(0,\hLoad)$) {$p$};
				
				%% ── Coppia antioraria in F ───────────────────
				\def\rcop{0.32}
				\draw[->, >=stealth, line width=0.9pt]
				($(F)+(\rcop,0)$) arc(0:270:\rcop);
				\node[right, font=\small]
				at ($(F)+(\rcop+0.06,0.08)$) {$p\ell^2$};
				
				%% ── Reazioni in B ────────────────────────────
				%% H_B = (7/6)pℓ  verso sinistra (←)
				\draw[reaz] (B) -- ++(-\sc, 0)
				node[above, font=\small, gray!70]
				{$\tfrac{7}{6}p\ell$};
				%% V_B = (2/3)pℓ  verso l'alto (↑)
				\draw[reaz] (B) -- ++(0, \sc)
				node[right, font=\small, gray!70]
				{$\tfrac{2}{3}p\ell$};
				
				%% ── Etichette nodi ───────────────────────────
				\node[above left,  font=\small] at (Ddes) {$D$};
				\node[above right, font=\small]
				at ($(E)+(0.05,0.08)$) {$E$};
				\node[right, font=\small] at (B)    {$B$};
				\node[left,        font=\small]
				at ($(F)+(-0.22,0)$)   {$F$};
				
			\end{scope}
			
		\end{tikzpicture}%
	}%
	\caption{Portale a tre cerniere: isolamento dei due semiarchi
		come corpi liberi.
		Frecce grigie: reazioni vincolari esterne.
		Frecce nere alle cerniere~$D$: azioni interne,
		opposte sui due semiarchi per il principio di
		azione e reazione.}
	\label{fig:arco3c_corpi_liberi}
\end{figure}
